%% Springer Nature LaTeX template (sn-jnl.cls). Build with pdflatex.
%% Reference style: sn-mathphys, numbered citations.
\documentclass[pdflatex,sn-mathphys,Numbered]{sn-jnl}

%%%% Additional packages used by this report
\usepackage[T1]{fontenc}
\usepackage[polish]{babel}
\usepackage{amsmath,amssymb}% sn-jnl.cls expects amsmath (\allowdisplaybreaks)
\usepackage{manyfoot}% required by sn-jnl.cls footnote machinery
\usepackage{booktabs}% sn-jnl.cls defines \toprule/\midrule but not \bottomrule
\usepackage{float}
\usepackage{enumitem}
\usepackage{tikz}
\usetikzlibrary{arrows.meta,positioning}
\usepackage{pgfplots}
\pgfplotsset{compat=1.18}
\usepgfplotslibrary{groupplots,statistics}

%% graphicx, amsmath, amssymb, xcolor, hyperref, natbib and the A4 page
%% geometry are already provided by sn-jnl.cls.

% Figure colours and the shared pgfplots `wykres' style live with the figures.
% Shared figure styling for the report. \input from the main preamble, AFTER
% xcolor/tikz/pgfplots are loaded. Both the report and the figure fragments in
% this directory rely on the colours and the `wykres` style defined here.

\definecolor{colOX}{HTML}{1F77B4}
\definecolor{colPMX}{HTML}{FF7F0E}
\definecolor{colOXC}{HTML}{2CA02C}
\definecolor{colPMXC}{HTML}{D62728}

\pgfplotsset{
    wykres/.style={
        width=0.92\textwidth, height=6.2cm,
        tick label style={font=\small}, label style={font=\small},
        legend style={font=\small, draw=none, fill=none},
        grid=major, grid style={gray!18},
        every axis plot/.append style={line width=1pt},
    },
}


% Lightweight, dependency-free pseudocode: a custom "Algorytm" float (float pkg)
% plus a numbered-line environment (enumitem). Avoids algorithm/algorithmicx,
% which are absent from minimal TeX installs, while still rendering inline math.
\newfloat{algorithm}{htbp}{loa}
\floatname{algorithm}{Algorytm}
\newenvironment{pseudo}
  {\begin{enumerate}[leftmargin=2.6em, itemsep=1pt, topsep=2pt,
     label=\scriptsize\sffamily\arabic*]\small}
  {\end{enumerate}}
\newcommand{\Kw}[1]{\textbf{#1}}% pseudocode keyword
\newcommand{\In}[1]{\textbf{Wejście:} #1}
\newcommand{\Out}[1]{\textbf{Wyjście:} #1}
\newcommand{\ind}{\hspace*{1.5em}}% one indentation level

\raggedbottom

\begin{document}

\title[Memetyczny algorytm harmonogramowania zadań]{Harmonogramowanie zadań w~środowisku wieloprocesorowym z~ograniczeniem pierwszeństwa: algorytm memetyczny z~lokalną optymalizacją ścieżki krytycznej}

\author[1]{\fnm{Wojciech} \sur{Jasiński}}

\author[1]{\fnm{Piotr} \sur{Iśtok}}

\affil[1]{\orgname{Politechnika Warszawska}, \orgaddress{\city{Warszawa}, \country{Polska}}}

\maketitle

\sloppy

% =====================================================================
\section{Wprowadzenie i opis problemu}
% =====================================================================
Celem projektu jest opracowanie i~implementacja algorytmu ewolucyjnego (AE) dla
problemu harmonogramowania zadań w~środowisku wieloprocesorowym z~ograniczeniem
pierwszeństwa (\textit{ang. Task Scheduling Problem with Precedence Constraints}).
Zadanie należy do klasy problemów NP-trudnych~\cite{ullman1975}, co uzasadnia
wykorzystanie metaheurystyk.

\paragraph{Definicja formalna.}
Problem zdefiniowano jako czwórkę $(T, p, l, M)$, gdzie:
\begin{itemize}[itemsep=3pt, topsep=3pt]
    \item $T = \{t_1, t_2, \dots, t_n\}$ jest zbiorem $n$ zadań do wykonania;
    \item $p : T \times T \to \{0, 1\}$ koduje relację pierwszeństwa: $p(t_i, t_j) = 1$ oznacza, że zadanie $t_i$ musi zostać zakończone przed rozpoczęciem $t_j$. Zakładamy, że graf skierowany indukowany przez~$p$ jest acykliczny (DAG, \textit{ang. Directed Acyclic Graph}); w~przeciwnym razie zbiór rozwiązań dopuszczalnych jest pusty;
    \item $l: T \to \mathbb{N}$ przypisuje każdemu zadaniu czas trwania;
    \item $M \in \mathbb{N}$ jest stałą liczbą procesorów (grup), między które rozdzielamy zadania. Procesor wykonuje jedno zadanie naraz; nazwę „procesor'' stosujemy dla zgodności z~literaturą.
\end{itemize}

\noindent
Wyzwaniem jest taki podział zadań między $M$ procesorów, by zachować wszystkie
więzy pierwszeństwa i~zminimalizować czas zakończenia ostatniego zadania, czyli
\textit{makespan} ($C_{max}$).

\paragraph{Pojęcia harmonogramowe.}
Ponieważ procesory są nierozróżnialne, harmonogram dopuszczalny nie zależy tylko
od~$p$, lecz także od~tego, które zadania trafiają na ten sam procesor. Dwa
zadania wykonywane kolejno na tym samym procesorze łączy \emph{krawędź maszynowa}
-- nie jest to więz pierwszeństwa, lecz ograniczenie zasobowe wynikające z~tego,
że procesor jest zajęty. Ciąg takich krawędzi tworzy \emph{ścieżkę maszynową}.
O~wartości $C_{max}$ może decydować zarówno ścieżka pierwszeństwa, jak i~ścieżka
maszynowa; rozróżnienie to jest kluczowe dla zaproponowanej metody lokalnej.
Przykład na rysunku~\ref{fig:przyklad} pokazuje graf zależności i~jego harmonogram.

\begin{figure}[H]
\centering
% Editorial schematic (NOT data-driven): example dependency graph and its
% decoded Gantt schedule. Hand-authored; not produced by report_figures.py.
\begin{tikzpicture}[
    task/.style={circle, draw, minimum size=9mm, inner sep=0pt, font=\small, align=center},
    >={Stealth[length=2mm]}, edge/.style={->, draw=gray}, crit/.style={->, draw=colPMXC, very thick, dashed}]
  \node[task] (A) at (0,1.0)   {A\\[-2pt]\tiny 2};
  \node[task] (C) at (2.6,1.0) {C\\[-2pt]\tiny 3};
  \node[task] (B) at (0,-1.0)  {B\\[-2pt]\tiny 4};
  \node[task] (D) at (2.6,-1.0){D\\[-2pt]\tiny 4};
  \draw[->, draw=gray, thick] (A)--(C);
  \draw[crit] (B)--(C) node[midway, above left=-2pt, font=\tiny, text=colPMXC]{zasoby};
\end{tikzpicture}
\hspace{0.7cm}
\begin{tikzpicture}[font=\small]
  \def\gh{0.8}
  \draw[->] (0,-0.25) -- (7.6,-0.25) node[right]{$t$};
  \foreach \t in {0,...,7} \draw (\t,-0.2)--(\t,-0.3) node[below=-1pt]{\tiny\t};
  \node[left] at (0,1.0+\gh/2) {\footnotesize P1};
  \node[left] at (0,0.0+\gh/2) {\footnotesize P2};
  \draw[fill=colOX!20] (0,1.0) rectangle (2,1.0+\gh) node[midway]{A};
  \draw[fill=colOX!20] (2,1.0) rectangle (6,1.0+\gh) node[midway]{D};
  \draw[fill=colPMXC!30] (0,0.0) rectangle (4,0.0+\gh) node[midway]{\bfseries B};
  \draw[fill=colPMXC!30] (4,0.0) rectangle (7,0.0+\gh) node[midway]{\bfseries C};
  \draw[dashed, colPMXC] (7,-0.25) -- (7,2.0) node[above]{\footnotesize $C_{max}=7$};
\end{tikzpicture}

\caption{Przykład ($M=2$; liczby w~węzłach to czasy trwania). Po lewej graf
zależności (przerywana linia $B\!\to\!C$ to krawędź maszynowa, nie więz
pierwszeństwa); po prawej harmonogram permutacji \texttt{[A,B,D,C]}. Start C
opóźniony do $t=4$ przez zajętość P2, więc o~$C_{max}$ decyduje ścieżka maszynowa
$B\!\to\!C$. Dlatego ścieżkę krytyczną (metoda ścieżki krytycznej, CPM --
\textit{ang. Critical Path Method}) liczymy na harmonogramie, a~nie na samym grafie.}
\label{fig:przyklad}
\end{figure}

% =====================================================================
\section{Metoda}
\label{sec:metoda}
% =====================================================================
Zastosowano algorytm memetyczny~\cite{moscato1989}, który łączy globalne
przeszukiwanie przestrzeni rozwiązań typowe dla algorytmów genetycznych (GA)
z~optymalizacją lokalną opartą o~analizę ścieżki krytycznej.

\paragraph{Reprezentacja osobnika.}
Osobnik koduje kolejność wykonywania zadań: chromosomem jest permutacja
o~długości $|T|$, zawierająca wszystkie zadania w~kolejności priorytetu
umieszczania w~harmonogramie. Reprezentacja nie zawiera bezpośredniego przydziału
do procesorów -- wynika on z~działania dekodera. Każda permutacja zgodna
z~porządkiem topologicznym koduje rozwiązanie dopuszczalne.

\paragraph{Dekoder (\textit{list scheduling}).}
Zamianę chromosomu na harmonogram i~wyznaczenie $C_{max}$ wykonuje zachłanny
dekoder typu \textit{list scheduling}~\cite{graham1969}
(Algorytm~\ref{alg:decoder}). Kolejne zadania z~permutacji przydziela się
procesorowi zwalniającemu się najwcześniej, nie wcześniej niż zakończą się
wszyscy poprzednicy w~relacji~$p$. Wyliczone $C_{max}$ jest wartością funkcji
przystosowania; każde jego wyznaczenie to jedna \emph{ewaluacja funkcji celu}.

\begin{algorithm}[H]
\caption{Dekoder list scheduling}
\label{alg:decoder}
\In{permutacja $\pi$ (zgodna z~porządkiem topologicznym), $M$ procesorów, czasy $l$.}\\
\Out{czasy startu/końca i~$C_{max}$.}
\begin{pseudo}
\item $H \gets$ kopiec procesorów, każdy z~czasem zwolnienia $0$
\item \Kw{dla} $t \in \pi$ (w~kolejności permutacji):
\item \ind $f_p \gets \max\{\text{koniec}(u) : p(u,t)=1\}$ \quad($0$, gdy brak poprzedników)
\item \ind $(c, m) \gets$ wyjmij z~$H$ procesor o~najmniejszym czasie zwolnienia $c$
\item \ind $\text{start}(t) \gets \max(c, f_p)$;\quad $\text{koniec}(t) \gets \text{start}(t) + l(t)$
\item \ind zapamiętaj poprzednika maszynowego $t$ (ostatnie zadanie na~$m$); wstaw $(\text{koniec}(t), m)$ do~$H$
\item \Kw{zwróć} $C_{max} \gets \max_t \text{koniec}(t)$
\end{pseudo}
\end{algorithm}

\paragraph{Naprawa dopuszczalności (\textit{ready-list repair}).}
Krzyżowanie i~mutacja mogą wygenerować permutacje niezgodne z~porządkiem
topologicznym. Zamiast karać takie rozwiązania, naprawiamy je
(Algorytm~\ref{alg:repair}): przebudowujemy chromosom, traktując jego kolejność
genów jako \emph{listę priorytetów}. W~każdym kroku spośród zadań gotowych
(o~zaspokojonych poprzednikach) wybieramy to o~najmniejszym indeksie w~oryginalnej
permutacji. Wynik jest najbliższym porządkiem topologicznym zachowującym względny
priorytet pierwotnej permutacji, więc cała populacja pozostaje dopuszczalna.

\begin{algorithm}[H]
\caption{Naprawa ready-list (przebudowa permutacji do porządku topologicznego)}
\label{alg:repair}
\In{permutacja $\pi$ (potencjalnie niedopuszczalna), DAG $G$.}\\
\Out{permutacja $\pi'$ zgodna z~porządkiem topologicznym.}
\begin{pseudo}
\item $\mathrm{poz}(v) \gets$ indeks $v$ w~$\pi$ dla każdego $v$ \quad(klucz priorytetu)
\item $\mathrm{in}(v) \gets$ stopień wejściowy $v$ w~$G$;\quad $\pi' \gets [\,]$
\item $R \gets$ kopiec minimowy zadań z~$\mathrm{in}(v)=0$, kluczowany przez $\mathrm{poz}(v)$
\item \Kw{dopóki} $R \neq \emptyset$:
\item \ind $v \gets$ wyjmij z~$R$ zadanie o~najmniejszym $\mathrm{poz}(v)$; dopisz $v$ do~$\pi'$
\item \ind \Kw{dla} każdego następnika $w$ zadania $v$:
\item \ind\ind $\mathrm{in}(w) \gets \mathrm{in}(w)-1$;\quad \Kw{jeśli} $\mathrm{in}(w)=0$ \Kw{to} wstaw $w$ do~$R$ z~kluczem $\mathrm{poz}(w)$
\item \Kw{zwróć} $\pi'$
\end{pseudo}
\end{algorithm}

\paragraph{Operatory ewolucyjne.}
Część ewolucyjną zaimplementowano w~oparciu o~bibliotekę DEAP~\cite{deap2012}.
\emph{Selekcja} jest turniejowa~\cite{miller1995}. \emph{Krzyżowanie} to operatory
permutacyjne unikające dublowania zadań: porównujemy krzyżowanie porządkowe (OX,
\textit{ang. Order Crossover})~\cite{davis1985}, zachowujące względną kolejność,
z~krzyżowaniem z~częściowym mapowaniem (PMX, \textit{ang. Partially Mapped
Crossover})~\cite{goldberg1985}, opartym o~mieszanie pozycyjne. \emph{Mutacja} to
operator zamiany (\textit{swap}) -- losowa zamiana miejscami dwóch zadań.

\paragraph{Optymalizacja lokalna (CPM).}
Kluczowym elementem metody memetycznej jest moduł oparty o~analizę ścieżki
krytycznej~\cite{kelley1959} (Algorytm~\ref{alg:cpm}). W~odróżnieniu od~klasycznej
CPM liczonej na samym grafie zależności, stosujemy analizę \emph{świadomą
harmonogramu}: ścieżkę krytyczną wyznaczamy na faktycznie zdekodowanym
harmonogramie, śledząc na każdym kroku poprzednika wiążącego -- z~relacji
pierwszeństwa \emph{lub} z~zajętości procesora (krawędź maszynowa). Dzięki temu
analiza wykrywa również wąskie gardła zasobowe. \emph{Ruchome} zadania krytyczne
-- czyli te, które mają jeszcze zapas między pozycją ostatniego poprzednika
a~własną pozycją w~permutacji -- przesuwamy w~lewo (operacja \textit{pull-left});
jeśli poprawia to $C_{max}$, zmiana -- zgodnie z~ewolucją
Lamarckowską~\cite{whitley1994} -- zastępuje osobnika w~populacji.

\begin{algorithm}[H]
\caption{Lokalne przeszukiwanie CPM (pull-left zadań krytycznych)}
\label{alg:cpm}
\In{permutacja $\pi$, budżet prób $K$.}\\
\Out{permutacja po~zaakceptowanych ruchach.}
\begin{pseudo}
\item $\pi^\ast \gets \pi$;\quad $C^\ast \gets \textsc{Dekoduj}(\pi^\ast)$
\item $\mathit{krit} \gets$ ruchome zadania krytyczne harmonogramu $\pi^\ast$ \quad(ścieżka krytyczna świadoma harmonogramu)
\item \Kw{powtórz} $K$ razy, dopóki $\mathit{krit} \neq \emptyset$:
\item \ind $t \gets$ losowe zadanie z~$\mathit{krit}$;\quad $b \gets$ maks.\ indeks poprzedników $t$ w~$p$ (lewa granica)
\item \ind \Kw{jeśli} $b \ge \mathrm{poz}(t)-1$ \Kw{to} pomiń próbę
\item \ind $j \gets$ losowy indeks z~$[\,b+1,\ \mathrm{poz}(t)-1\,]$;\quad $\pi' \gets$ przesuń $t$ na~pozycję $j$
\item \ind \Kw{jeśli} $\textsc{Dekoduj}(\pi') \le C^\ast$ \Kw{to} $\pi^\ast \gets \pi'$;\ $C^\ast \gets \textsc{Dekoduj}(\pi')$;\ przelicz $\mathit{krit}$
\item \Kw{zwróć} $\pi^\ast$
\end{pseudo}
\end{algorithm}

\paragraph{Schemat algorytmu memetycznego.}
Algorytm~\ref{alg:memetic} spina powyższe komponenty. Z~połączenia dwóch
operatorów krzyżowania (OX/PMX) oraz obecności bądź braku modułu CPM powstają
cztery porównywane konfiguracje: \texttt{ox}, \texttt{pmx} (czysty AE) oraz
\texttt{ox\_cpm}, \texttt{pmx\_cpm} (pełny algorytm memetyczny).

\begin{algorithm}[H]
\caption{Pętla algorytmu memetycznego (warianty z~CPM i~bez)}
\label{alg:memetic}
\In{instancja, budżet ewaluacji $B$, operator krzyżowania, flaga CPM.}\\
\Out{najlepszy znaleziony osobnik.}
\begin{pseudo}
\item populacja $\gets$ losowe permutacje topologiczne; oceń populację
\item \Kw{dopóki} liczba ewaluacji $< B$:
\item \ind potomstwo $\gets$ selekcja turniejowa(populacja); krzyżowanie; mutacja
\item \ind \textsc{NaprawaReadyList}(potomstwo) \quad(Algorytm~\ref{alg:repair})
\item \ind \Kw{jeśli} wariant memetyczny \Kw{to} \textsc{LokalnePrzeszukiwanieCPM}(potomstwo) \quad(Algorytm~\ref{alg:cpm})
\item \ind oceń potomstwo; populacja $\gets$ potomstwo
\item \Kw{zwróć} najlepszy znaleziony osobnik
\end{pseudo}
\end{algorithm}

% =====================================================================
\section{Plan eksperymentów}
\label{sec:eksperymenty}
% =====================================================================
\paragraph{Cel i hipotezy badawcze.}
Eksperymenty oceniają jakość rozwiązań algorytmu oraz weryfikują dwie hipotezy:
\begin{itemize}[itemsep=1pt, topsep=2pt, leftmargin=3em]
    \item[(H1)] lokalna optymalizacja CPM istotnie obniża medianę $C_{max}$ względem czystego AE \emph{przy równym budżecie ewaluacji funkcji celu};
    \item[(H2)] krzyżowanie OX istotnie dominuje nad PMX na grafach rzadkich ($|E|/|T|^2 < 0{,}1$).
\end{itemize}

\paragraph{Zestawy testowe.}
Jako instancje posłużyły grafy ze zbioru \textit{Standard Task Graph Set}
(Kasahara Lab)~\cite{tobita2002} oraz grafy generowane syntetycznie.
Generator tworzy losowy DAG: $n$ zadań o~czasach z~rozkładu jednostajnego
$[1,10]$, krawędzie losowane spośród par $(i,j)$ z~$i<j$ (co gwarantuje
acykliczność) w~liczbie zadanej gęstością $|E|/\binom{n}{2}$, a~węzły bez
poprzedników/następników spinane są sztucznymi węzłami wejścia i~wyjścia
o~czasie~$0$; ustalone ziarno czyni zbiór powtarzalnym. Dla małych instancji (do~$32$ węzłów
łącznie ze sztucznymi węzłami wejścia/wyjścia) wyznaczono rozwiązanie optymalne
solverem programowania całkowitoliczbowego mieszanego (MILP, \textit{ang.
Mixed-Integer Linear Programming}; CBC poprzez PuLP~\cite{mitchell2011}), co
pozwala raportować lukę optymalizacyjną. Dla pozostałych instancji odniesieniem
jest dolne ograniczenie (LB, \textit{ang. lower bound})
$\mathrm{LB} = \max(\lceil \text{całkowita praca}/M \rceil,\ \text{długość ścieżki krytycznej})$.
Luki dla różnych instancji nie są więc wprost porównywalne (raz względem optimum,
raz względem ograniczenia o~zmiennej szczelności); używamy ich do porównań
\emph{w~obrębie} instancji.

\paragraph{Budżet funkcji celu (kluczowe).}
Warianty memetyczne wykonują dodatkowe ewaluacje funkcji celu wewnątrz
przeszukiwania CPM (każda próba \textit{pull-left} to ponowne dekodowanie), więc
przy równej liczbie \emph{pokoleń} dysponowałyby wielokrotnie większym nakładem
obliczeniowym. Aby porównanie było uczciwe~\cite{hooker1995}, wszystkie cztery
konfiguracje prowadzimy do tego samego \emph{budżetu ewaluacji funkcji celu}
$B = 5050$ dekodowań harmonogramu (tyle zużywa czysty AE przy populacji $50$
i~$100$ pokoleniach: $50\times(100+1)$). Każdy przebieg zatrzymujemy na granicy
pokolenia po~osiągnięciu~$B$. Dzięki temu wariant memetyczny wykonuje mniej
pokoleń, ale w~przybliżeniu tyle samo ewaluacji co~AE -- różnice jakości nie
wynikają już z~większego nakładu. Ponieważ jedno pokolenie wariantu memetycznego
to wiele ewaluacji, warianty te nieznacznie przekraczają budżet (średnio ok.~5650
wobec 5070 ewaluacji czystego AE) -- dostają więc \emph{nieco więcej}, a~nie mniej
nakładu, co czyni wniosek o~przewadze czystego AE konserwatywnym. Dla porównania
przytaczamy też wcześniejsze wyniki przy równej liczbie pokoleń.

\paragraph{Protokół statystyczny.}
Dla każdej konfiguracji wykonano 15 niezależnych uruchomień (ziarna $0$--$14$);
miarą jakości jest najlepszy $C_{max}$ na uruchomienie, raportujemy medianę.
Istotność różnic weryfikowano sparowanym testem kolejności rangowanej
Wilcoxona~\cite{wilcoxon1945} (implementacja \texttt{scipy}~\cite{virtanen2020};
te same ziarna w~obu porównywanych konfiguracjach), na~poziomie $\alpha=0{,}05$.
Raportujemy surową p-wartość każdego testu \emph{osobno dla instancji} i~nie
stosujemy korekcji na~wielokrotne porównania -- przy 20 instancjach pojedyncze
fałszywie istotne wyniki przy $\alpha=0{,}05$ są możliwe, dlatego wnioski opieramy
na~rozmiarze efektu i~spójnym kierunku różnic, a~nie na~pojedynczych wynikach
istotności. Rozmiar efektu podajemy jako różnicę median $\Delta$, bo przy 15
ziarnach nawet różnica jednej jednostki bywa istotna statystycznie, choć
praktycznie znikoma. Punktem odniesienia są heurystyki listowe HLFET (\textit{ang.
Highest Level First with Estimated Times})~\cite{adam1974} i~MCP (\textit{ang.
Modified Critical Path})~\cite{wu1990} oraz czysty AE.

\paragraph{Hiperparametry.}
Bazowe parametry AE są globalne (jednakowe dla wszystkich instancji); ich wartości
i~uzasadnienie -- zakresy ugruntowane w~literaturze GA oraz prac
o~harmonogramowaniu~\cite{dejong1990, grefenstette1986, omara2010, hou1994} --
zestawia tabela~\ref{tab:hparams}. Te wartości \emph{poddajemy następnie
strojeniu} (sekcja~\ref{sec:wyniki}).

\begin{table}[ht]
\centering\footnotesize\setlength{\tabcolsep}{6pt}
\begin{tabular}{llp{8.0cm}}
\toprule
parametr & wartość & podstawa \\
\midrule
prawd. krzyżowania $p_c$ & $0{,}8$ & klasyczne $0{,}6$~\cite{dejong1990}--$0{,}9$~\cite{grefenstette1986}; GA dla DAG $0{,}65$--$0{,}95$~\cite{omara2010, hou1994} \\
prawd. mutacji $p_m$ & $0{,}3$ & zamiana (\textit{swap}) na osobnika; GA permutacyjne $0{,}1$--$0{,}5$ \\
populacja & $50$ & \cite{dejong1990}; powszechny zakres $30$--$100$ \\
budżet ewaluacji $B$ & $5050$ & $50\times(100+1)$ -- nakład czystego AE (wspólny dla wszystkich konfiguracji) \\
rozmiar turnieju & $3$ & łagodny nacisk selekcyjny~\cite{miller1995} \\
próby CPM / osobnika & $20$ & budżet lokalnego przeszukiwania (warianty memetyczne) \\
niezależne ziarna & $15$ & ograniczone czasem obliczeń; poniżej konwencji $30$ \\
\bottomrule
\end{tabular}
\caption{Bazowe hiperparametry algorytmu i~ich uzasadnienie. Wartości są globalne
dla wszystkich instancji i~stanowią punkt wyjścia dla strojenia.}
\label{tab:hparams}
\end{table}

% =====================================================================
\section{Wyniki}
\label{sec:wyniki}
% =====================================================================
\paragraph{Jakość rozwiązań przy równym budżecie.}
Tabela~\ref{tab:main} zbiera medianę $C_{max}$ dla czterech konfiguracji przy
równym budżecie ewaluacji. Przy uczciwym porównaniu \emph{najlepszy jest czysty
\texttt{pmx}}: jest najlepszy lub remisuje na 19 z~20 instancji (przegrywa tylko
o~$1$ na \texttt{syn\_t22} z~\texttt{pmx\_cpm}) oraz dorównuje lub pokonuje
heurystyki HLFET/MCP na 19 z~20 (wyjątkiem jest \texttt{robot\_t88}, gdzie
heurystyka $694$ bije wszystkie warianty AE). Część instancji jest wymuszona --
każda konfiguracja osiąga ograniczenie (np. \texttt{rand0014\_t500}: wszyscy $893$).

\begin{table}[ht]
\centering\scriptsize\setlength{\tabcolsep}{3pt}
\begin{tabular}{lrrrrrrrrr}
\toprule
instancja & $|T|$ & gęstość & $M$ & ref. & HLFET/MCP & \texttt{ox} & \texttt{pmx} & \texttt{ox\_cpm} & \texttt{pmx\_cpm} \\
\midrule
\texttt{syn\_t15\_p010}       &  17 & 0{,}093 & 2 &  52 &   54 & \textbf{52}  & \textbf{52}  & \textbf{52}  & \textbf{52} \\
\texttt{syn\_t22\_p018}       &  24 & 0{,}090 & 3 &  44 &   48 & 46           & 45           & 45           & \textbf{44} \\
\texttt{syn\_t30\_p030}       &  32 & 0{,}132 & 3 &  66 &   66 & \textbf{66}  & \textbf{66}  & \textbf{66}  & \textbf{66} \\
\texttt{syn\_t40\_p006}       &  42 & 0{,}044 & 4 &  52 &   55 & 55           & \textbf{52}  & 53           & \textbf{52} \\
\texttt{syn\_t40\_p035}       &  42 & 0{,}157 & 4 & 108 &  108 & \textbf{108} & \textbf{108} & \textbf{108} & \textbf{108} \\
\texttt{syn\_t60\_p005}       &  62 & 0{,}033 & 4 &  72 &   74 & 81           & \textbf{73}  & 77           & 75 \\
\texttt{syn\_t60\_p015}       &  62 & 0{,}073 & 4 &  77 &   88 & 89           & \textbf{81}  & 84           & 82 \\
\texttt{rand0000\_t50\_p008}  &  52 & 0{,}044 & 4 &  66 &   77 & 73           & \textbf{67}  & 69           & 68 \\
\texttt{rand0093\_t50\_p078}  &  52 & 0{,}360 & 4 &  69 &   73 & \textbf{73}  & \textbf{73}  & \textbf{73}  & \textbf{73} \\
\texttt{rand0063\_t100\_p002} & 102 & 0{,}017 & 4 & 136 &  137 & 137          & \textbf{136} & \textbf{136} & \textbf{136} \\
\texttt{rand0173\_t100\_p039} & 102 & 0{,}186 & 4 & 199 &  218 & 227          & \textbf{213} & 224          & 220 \\
\texttt{syn\_t100\_p004}      & 102 & 0{,}024 & 4 & 126 &  126 & 128          & \textbf{126} & 127          & \textbf{126} \\
\texttt{syn\_t100\_p012}      & 102 & 0{,}059 & 4 & 131 &  153 & 162          & \textbf{142} & 150          & 147 \\
\texttt{syn\_t150\_p004}      & 152 & 0{,}021 & 4 & 230 &  241 & 234          & \textbf{230} & 233          & 231 \\
\texttt{rand0071\_t300\_p001} & 302 & 0{,}006 & 4 & 423 &  428 & 425          & \textbf{423} & 424          & \textbf{423} \\
\texttt{rand0000\_t300\_p016} & 302 & 0{,}080 & 4 & 448 &  505 & 531          & \textbf{502} & 523          & 512 \\
\texttt{robot\_t88\_p003}     &  90 & 0{,}017 & 4 & 621 &  694 & 814          & \textbf{741} & 805          & 802 \\
\texttt{sparse\_t96\_p001}    &  98 & 0{,}018 & 4 & 484 &  486 & 502          & \textbf{486} & 494          & 490 \\
\texttt{fpppp\_t334\_p002}    & 336 & 0{,}011 & 4 &1783 & 1903 & 2067         & \textbf{1894}& 2064         & 2036 \\
\texttt{rand0014\_t500\_p020} & 502 & 0{,}098 & 8 & 893 &  893 & \textbf{893} & \textbf{893} & \textbf{893} & \textbf{893} \\
\bottomrule
\end{tabular}
\caption{Mediana $C_{max}$ wg~konfiguracji \textbf{przy równym budżecie ewaluacji}
($B=5050$; pogrubienie -- najlepsza na instancję). $M$ -- liczba procesorów.
Odniesienie: optimum MILP dla instancji do~$32$ węzłów, w~pozostałych przypadkach
dolne ograniczenie LB. Kolumna HLFET/MCP podaje wynik heurystyk listowych -- obie
pokrywają się na każdej instancji.}
\label{tab:main}
\end{table}

\paragraph{Dlaczego heurystyki HLFET i~MCP dają identyczne wyniki.}
Obie heurystyki różnią się jedynie regułą priorytetu (poziom statyczny vs~najpóźniejszy
czas startu), lecz obie produkują pojedynczy porządek topologiczny podawany temu
samemu dekoderowi list scheduling. Na badanych instancjach z~nierozróżnialnymi
procesorami oba porządki prowadzą do~tego samego $C_{max}$, więc kolumny HLFET
i~MCP są identyczne; nie jest to błąd, lecz efekt dekodera niewrażliwego na~te
akurat różnice priorytetów.

\paragraph{H1 -- wpływ CPM przy równym budżecie.}
Hipoteza H1 \emph{nie potwierdza się} w~uczciwym porównaniu (tabela~\ref{tab:h1}).
Na 40 testach (2~operatory $\times$ 20 instancji) istotnych jest 25: CPM pomaga
w~16, a~szkodzi w~9. Podział jest jednak wyraźnie operatorowy. Dla \emph{słabszego}
OX CPM pomaga na~wszystkich 15 nieremisowych instancjach (0 pogorszeń), bo lokalne
przeszukiwanie nadrabia braki eksploracji. Dla \emph{silniejszego} PMX --
przeciwnie: istotnie szkodzi na~9 instancjach, a~pomaga tylko raz (o~1 jednostkę na
\texttt{syn\_t22}). Praktycznie istotne ($|\Delta|\ge 5$) jest pięć pogorszeń PMX
-- \texttt{fpppp} ($\Delta=-142$), \texttt{robot\_t88} ($-61$),
\texttt{rand0000\_t300} ($-10$), \texttt{rand0173} ($-7$), \texttt{syn\_t100\_p012}
($-5$) -- pozostałe to różnice $1$--$2$ jednostek. Zbieżność zbiorcza w~funkcji
liczby ewaluacji (rys.~\ref{fig:convevals}) pokazuje mechanizm: przy tym samym
budżecie czysty \texttt{pmx} schodzi najniżej, bo przeznacza wszystkie ewaluacje
na~eksplorację wielu pokoleń, podczas gdy \texttt{pmx\_cpm} zużywa budżet
na~intensyfikację w~niewielu pokoleniach.

\begin{table}[ht]
\centering\scriptsize\setlength{\tabcolsep}{5pt}
\begin{tabular}{lrrrr}
\toprule
& \multicolumn{2}{c}{OX} & \multicolumn{2}{c}{PMX} \\
\cmidrule(lr){2-3}\cmidrule(lr){4-5}
instancja & $\Delta$ & $p$ & $\Delta$ & $p$ \\
\midrule
\texttt{syn\_t15\_p010} & +0 & 1{,}000 & +0 & 1{,}000 \\
\texttt{syn\_t22\_p018} & \textbf{+1} & \textbf{0{,}001} & \textbf{+1} & \textbf{0{,}005} \\
\texttt{syn\_t30\_p030} & +0 & 1{,}000 & +0 & 1{,}000 \\
\texttt{syn\_t40\_p006} & \textbf{+2} & \textbf{0{,}001} & +0 & 0{,}564 \\
\texttt{syn\_t40\_p035} & +0 & 1{,}000 & +0 & 1{,}000 \\
\texttt{syn\_t60\_p005} & \textbf{+4} & \textbf{0{,}001} & \textbf{$-2$} & \textbf{0{,}001} \\
\texttt{syn\_t60\_p015} & \textbf{+5} & \textbf{0{,}001} & \textbf{$-1$} & \textbf{0{,}001} \\
\texttt{rand0000\_t50\_p008} & \textbf{+4} & \textbf{0{,}001} & \textbf{$-1$} & \textbf{0{,}001} \\
\texttt{rand0093\_t50\_p078} & +0 & 1{,}000 & +0 & 1{,}000 \\
\texttt{rand0063\_t100\_p002} & \textbf{+1} & \textbf{0{,}008} & +0 & 0{,}317 \\
\texttt{rand0173\_t100\_p039} & \textbf{+3} & \textbf{0{,}001} & \textbf{$-7$} & \textbf{0{,}001} \\
\texttt{syn\_t100\_p004} & \textbf{+1} & \textbf{0{,}001} & +0 & 1{,}000 \\
\texttt{syn\_t100\_p012} & \textbf{+12} & \textbf{0{,}001} & \textbf{$-5$} & \textbf{0{,}002} \\
\texttt{syn\_t150\_p004} & \textbf{+1} & \textbf{0{,}004} & \textbf{$-1$} & \textbf{0{,}005} \\
\texttt{rand0071\_t300\_p001} & \textbf{+1} & \textbf{0{,}004} & +0 & 0{,}058 \\
\texttt{rand0000\_t300\_p016} & \textbf{+8} & \textbf{0{,}001} & \textbf{$-10$} & \textbf{0{,}001} \\
\texttt{robot\_t88\_p003} & \textbf{+9} & \textbf{0{,}011} & \textbf{$-61$} & \textbf{0{,}001} \\
\texttt{sparse\_t96\_p001} & \textbf{+8} & \textbf{0{,}001} & $-4$ & 0{,}182 \\
\texttt{fpppp\_t334\_p002} & \textbf{+3} & \textbf{0{,}019} & \textbf{$-142$} & \textbf{0{,}001} \\
\texttt{rand0014\_t500\_p020} & +0 & 1{,}000 & +0 & 1{,}000 \\
\bottomrule
\end{tabular}
\caption{H1 -- test Wilcoxona CPM-on vs CPM-off, osobno dla~OX i~PMX.
$\Delta = \text{mediana}_{\text{off}} - \text{mediana}_{\text{on}}$ (dodatnie =
CPM pomaga); $p$ -- surowa p-wartość Wilcoxona. Pogrubienie -- istotne
($p<0{,}05$). Dla~OX wszystkie istotne różnice są dodatnie (CPM pomaga), dla~PMX
prawie wszystkie ujemne (CPM szkodzi).}
\label{tab:h1}
\end{table}

\begin{figure}[H]
\centering
% Generated by src/experiments/report_figures.py — do not edit by hand.
% Regenerate: python -m src.experiments.report_figures --campaign results/<dir>
\begin{tikzpicture}
\begin{axis}[wykres, height=5.6cm, xlabel={ewaluacje funkcji celu},
    ylabel={średnia luka do LB [\%]},
    xmin=50, xmax=5050, ymin=0, legend pos=north east,
    legend style={font=\footnotesize}]
\addplot[colOX, thick] table[x=evals, y=ox] {figures/data/conv_evals_aggregate.dat};
\addplot[colPMX, thick] table[x=evals, y=pmx] {figures/data/conv_evals_aggregate.dat};
\addplot[colOXC, thick] table[x=evals, y=ox_cpm] {figures/data/conv_evals_aggregate.dat};
\addplot[colPMXC, thick] table[x=evals, y=pmx_cpm] {figures/data/conv_evals_aggregate.dat};
\legend{ox,pmx,ox\_cpm,pmx\_cpm}
\end{axis}
\end{tikzpicture}

\caption{Zbieżność zbiorcza przy równym budżecie: średnia luka do~LB po~wszystkich
20 instancjach w~funkcji liczby ewaluacji funkcji celu. Przy tym samym budżecie
czysty \texttt{pmx} kończy najniżej; CPM nie nadrabia różnicy na korzyść PMX.}
\label{fig:convevals}
\end{figure}

\paragraph{Artefakt budżetowy.}
Odwrócenie H1 jest skutkiem nierównego budżetu we~wcześniejszej wersji
eksperymentu. Przy równej liczbie pokoleń (po~$100$) warianty memetyczne wykonywały
ok.~$20$--$40\times$ więcej ewaluacji na~pokolenie i~wtedy \texttt{pmx\_cpm}
wygrywał wszędzie (np.~\texttt{fpppp} $1843$, \texttt{robot} $697$,
\texttt{syn\_t100\_p012} $135$). Po~wyrównaniu budżetu te same instancje dają
\texttt{pmx\_cpm} $2036$, $802$, $147$ -- gorzej niż czysty \texttt{pmx}
($1894$, $741$, $142$). Pozorna przewaga CPM była więc artefaktem większego nakładu,
a~nie lepszej metody.

\paragraph{H2 -- OX vs PMX.}
Hipoteza H2 również się nie potwierdza i~zachodzi relacja przeciwna, niezależnie
od~budżetu (tabela~\ref{tab:h2}). PMX istotnie pokonuje OX na~każdej nieremisowej
instancji: na~16 grafach rzadkich -- do~których odnosi się hipoteza -- to 14
istotnych zwycięstw PMX, 2 remisy i~0 dla~OX; na~wszystkich 20 instancjach -- 15
istotnych zwycięstw, 5 remisów, 0 dla~OX. Remisy występują tam, gdzie oba operatory
osiągają ograniczenie. Prawdopodobny mechanizm: naprawa listą gotowych
(Algorytm~\ref{alg:repair}) już wymusza dopuszczalny porządek, więc atut OX
w~zachowaniu kolejności jest stłumiony, podczas gdy pozycyjne mieszanie PMX
skuteczniej przeszukuje przestrzeń priorytetów. Jest to hipoteza wyjaśniająca
obserwację, nie wynik osobnego eksperymentu ablacyjnego.

\begin{table}[ht]
\centering\scriptsize\setlength{\tabcolsep}{5pt}
\begin{tabular}{lrrrrr}
\toprule
instancja & gęstość & \texttt{ox} & \texttt{pmx} & $\Delta$ & $p$ \\
\midrule
\texttt{syn\_t15\_p010} & 0{,}093 & 52 & 52 & +0 & 1{,}000 \\
\texttt{syn\_t22\_p018} & 0{,}090 & 46 & 45 & \textbf{$-1$} & \textbf{0{,}009} \\
\texttt{syn\_t30\_p030} & 0{,}132 & 66 & 66 & +0 & 1{,}000 \\
\texttt{syn\_t40\_p006} & 0{,}044 & 55 & 52 & \textbf{$-3$} & \textbf{0{,}001} \\
\texttt{syn\_t40\_p035} & 0{,}157 & 108 & 108 & +0 & 1{,}000 \\
\texttt{syn\_t60\_p005} & 0{,}033 & 81 & 73 & \textbf{$-8$} & \textbf{0{,}001} \\
\texttt{syn\_t60\_p015} & 0{,}073 & 89 & 81 & \textbf{$-8$} & \textbf{0{,}001} \\
\texttt{rand0000\_t50\_p008} & 0{,}044 & 73 & 67 & \textbf{$-6$} & \textbf{0{,}001} \\
\texttt{rand0093\_t50\_p078} & 0{,}360 & 73 & 73 & +0 & 1{,}000 \\
\texttt{rand0063\_t100\_p002} & 0{,}017 & 137 & 136 & \textbf{$-1$} & \textbf{$<$0{,}001} \\
\texttt{rand0173\_t100\_p039} & 0{,}186 & 227 & 213 & \textbf{$-14$} & \textbf{0{,}001} \\
\texttt{syn\_t100\_p004} & 0{,}024 & 128 & 126 & \textbf{$-2$} & \textbf{0{,}001} \\
\texttt{syn\_t100\_p012} & 0{,}059 & 162 & 142 & \textbf{$-20$} & \textbf{0{,}001} \\
\texttt{syn\_t150\_p004} & 0{,}021 & 234 & 230 & \textbf{$-4$} & \textbf{0{,}001} \\
\texttt{rand0071\_t300\_p001} & 0{,}006 & 425 & 423 & \textbf{$-2$} & \textbf{0{,}001} \\
\texttt{rand0000\_t300\_p016} & 0{,}080 & 531 & 502 & \textbf{$-29$} & \textbf{0{,}001} \\
\texttt{robot\_t88\_p003} & 0{,}017 & 814 & 741 & \textbf{$-73$} & \textbf{0{,}001} \\
\texttt{sparse\_t96\_p001} & 0{,}018 & 502 & 486 & \textbf{$-16$} & \textbf{0{,}001} \\
\texttt{fpppp\_t334\_p002} & 0{,}011 & 2067 & 1894 & \textbf{$-173$} & \textbf{0{,}001} \\
\texttt{rand0014\_t500\_p020} & 0{,}098 & 893 & 893 & +0 & 1{,}000 \\
\bottomrule
\end{tabular}
\caption{H2 -- test Wilcoxona OX vs PMX (czysty AE, równy budżet).
$\Delta = \text{mediana}_{\text{pmx}} - \text{mediana}_{\text{ox}}$ (ujemne =
PMX lepszy); $p$ -- surowa p-wartość Wilcoxona. Pogrubienie -- istotne
($p<0{,}05$). PMX nie przegrywa z~OX na~żadnej instancji.}
\label{tab:h2}
\end{table}

\paragraph{Rozrzut i ujęcie zbiorcze.}
Boxploty (rys.~\ref{fig:box}) potwierdzają, że \texttt{pmx} jest nie tylko
najniższy, ale i~wąski. Rysunek~\ref{fig:agg} porządkuje konfiguracje średnią luką
do~LB: w~każdym podziale (wszystkie/gęste/rzadkie) czysty \texttt{pmx} jest
najlepszy ($3{,}5\%$ na~wszystkich instancjach wobec $5{,}1\%$ dla \texttt{pmx\_cpm},
$6{,}2\%$ dla \texttt{ox\_cpm} i~$8{,}3\%$ dla \texttt{ox}).

\begin{figure}[H]
\centering
% Generated by src/experiments/report_figures.py — do not edit by hand.
% Regenerate: python -m src.experiments.report_figures --campaign results/<dir>
\begin{tikzpicture}
\begin{groupplot}[
    group style={group size=2 by 2, vertical sep=1.7cm, horizontal sep=1.5cm},
    width=0.5\textwidth, height=4.4cm, boxplot/draw direction=y,
    xtick={1,2,3,4}, xticklabels={ox,pmx,ox\_cpm,pmx\_cpm},
    x tick label style={font=\footnotesize, rotate=25, anchor=east},
    tick label style={font=\footnotesize}, ylabel style={font=\small},
    ylabel={$C_{max}$},
    ymajorgrids=true, grid style={gray!18},
]
\nextgroupplot[title={\small \texttt{syn\_t100\_p012}}]
\addplot[draw=colOX,fill=colOX!25,boxplot prepared={lower whisker=158,lower quartile=161,median=162,upper quartile=163,upper whisker=164}]coordinates{};
\addplot[draw=colPMX,fill=colPMX!25,boxplot prepared={lower whisker=138,lower quartile=142,median=142,upper quartile=143,upper whisker=147}]coordinates{};
\addplot[draw=colOXC,fill=colOXC!25,boxplot prepared={lower whisker=144,lower quartile=150,median=150,upper quartile=153,upper whisker=155}]coordinates{};
\addplot[draw=colPMXC,fill=colPMXC!25,boxplot prepared={lower whisker=143,lower quartile=145.5,median=147,upper quartile=147,upper whisker=149}]coordinates{};
\nextgroupplot[title={\small \texttt{rand0173\_t100\_p039}}]
\addplot[draw=colOX,fill=colOX!25,boxplot prepared={lower whisker=224,lower quartile=226.5,median=227,upper quartile=228.5,upper whisker=231}]coordinates{};
\addplot[draw=colPMX,fill=colPMX!25,boxplot prepared={lower whisker=211,lower quartile=212.5,median=213,upper quartile=215,upper whisker=216}]coordinates{};
\addplot[draw=colOXC,fill=colOXC!25,boxplot prepared={lower whisker=220,lower quartile=222.5,median=224,upper quartile=225,upper whisker=228}]coordinates{};
\addplot[draw=colPMXC,fill=colPMXC!25,boxplot prepared={lower whisker=216,lower quartile=218,median=220,upper quartile=221,upper whisker=222}]coordinates{};
\nextgroupplot[title={\small \texttt{robot\_t88\_p003}}]
\addplot[draw=colOX,fill=colOX!25,boxplot prepared={lower whisker=784,lower quartile=810.5,median=814,upper quartile=823.5,upper whisker=834}]coordinates{};
\addplot[draw=colPMX,fill=colPMX!25,boxplot prepared={lower whisker=718,lower quartile=733,median=741,upper quartile=751,upper whisker=761}]coordinates{};
\addplot[draw=colOXC,fill=colOXC!25,boxplot prepared={lower whisker=780,lower quartile=795,median=805,upper quartile=811,upper whisker=821}]coordinates{};
\addplot[draw=colPMXC,fill=colPMXC!25,boxplot prepared={lower whisker=797,lower quartile=798.5,median=802,upper quartile=806,upper whisker=814}]coordinates{};
\nextgroupplot[title={\small \texttt{fpppp\_t334\_p002}}]
\addplot[draw=colOX,fill=colOX!25,boxplot prepared={lower whisker=2038,lower quartile=2060,median=2067,upper quartile=2079.5,upper whisker=2090}]coordinates{};
\addplot[draw=colPMX,fill=colPMX!25,boxplot prepared={lower whisker=1827,lower quartile=1859.5,median=1894,upper quartile=1923.5,upper whisker=1942}]coordinates{};
\addplot[draw=colOXC,fill=colOXC!25,boxplot prepared={lower whisker=2024,lower quartile=2040.5,median=2064,upper quartile=2068.5,upper whisker=2076}]coordinates{};
\addplot[draw=colPMXC,fill=colPMXC!25,boxplot prepared={lower whisker=1955,lower quartile=2028,median=2036,upper quartile=2046.5,upper whisker=2059}]coordinates{};
\end{groupplot}
\end{tikzpicture}

\caption{Rozrzut $C_{max}$ po~15 ziarnach (równy budżet). \texttt{pmx} -- najniższy.}
\label{fig:box}
\end{figure}

\begin{figure}[H]
\centering
% Generated by src/experiments/report_figures.py — do not edit by hand.
% Regenerate: python -m src.experiments.report_figures --campaign results/<dir>
\begin{tikzpicture}
\begin{groupplot}[group style={group size=3 by 1, horizontal sep=1.0cm},
    width=0.32\textwidth, height=4.16cm, ybar, ymin=0, ymax=10,
    symbolic x coords={ox,pmx,ox\_cpm,pmx\_cpm}, xtick={ox,pmx,ox\_cpm,pmx\_cpm}, enlarge x limits=0.18,
    x tick label style={rotate=25, anchor=east, font=\footnotesize},
    tick label style={font=\footnotesize},
    title style={font=\small}, label style={font=\small},
    ymajorgrids=true, grid style={gray!18},
    nodes near coords, every node near coord/.append style={font=\tiny},
    every axis plot/.append style={bar shift=0pt, bar width=9pt}]
\nextgroupplot[title={wszystkie (20)}, ylabel={średnia luka do LB [\%]}]
\addplot[fill=colOX,draw=colOX] coordinates {(ox,8.32)};
\addplot[fill=colPMX,draw=colPMX] coordinates {(pmx,3.48)};
\addplot[fill=colOXC,draw=colOXC] coordinates {(ox\_cpm,6.21)};
\addplot[fill=colPMXC,draw=colPMXC] coordinates {(pmx\_cpm,5.08)};
\nextgroupplot[title={gęste ($\ge 0{,}1$)}]
\addplot[fill=colOX,draw=colOX] coordinates {(ox,4.97)};
\addplot[fill=colPMX,draw=colPMX] coordinates {(pmx,3.21)};
\addplot[fill=colOXC,draw=colOXC] coordinates {(ox\_cpm,4.59)};
\addplot[fill=colPMXC,draw=colPMXC] coordinates {(pmx\_cpm,4.09)};
\nextgroupplot[title={rzadkie ($< 0{,}1$)}]
\addplot[fill=colOX,draw=colOX] coordinates {(ox,9.15)};
\addplot[fill=colPMX,draw=colPMX] coordinates {(pmx,3.55)};
\addplot[fill=colOXC,draw=colOXC] coordinates {(ox\_cpm,6.61)};
\addplot[fill=colPMXC,draw=colPMXC] coordinates {(pmx\_cpm,5.33)};
\end{groupplot}
\end{tikzpicture}

\caption{Średnia luka do~LB wg~konfiguracji (równy budżet): wszystkie instancje,
grafy gęste i~rzadkie. Czysty \texttt{pmx} jest najlepszy w~każdym podziale.}
\label{fig:agg}
\end{figure}

\paragraph{Koszt obliczeniowy.}
Przy równym budżecie koszt obliczeniowy czterech konfiguracji jest zbliżony
(rys.~\ref{fig:cpuscatter}) -- liczba dekodowań jest dominującym kosztem i~jest
wspólna. Znika tym samym wcześniejszy obraz, w~którym warianty memetyczne były
o~rzędy wielkości droższe: tamta różnica także była konsekwencją nierównego budżetu.

\begin{figure}[H]
\centering
% Generated by src/experiments/report_figures.py — do not edit by hand.
% Regenerate: python -m src.experiments.report_figures --campaign results/<dir>
\begin{tikzpicture}
\begin{axis}[
    width=0.92\textwidth, height=4.32cm, xmode=log, ymode=log,
    xlabel={liczba zadań $|T|$}, ylabel={średni czas CPU / uruchomienie [s]},
    tick label style={font=\small}, label style={font=\small},
    grid=major, grid style={gray!18},
    xmin=14, xmax=560, ymin=0.05, ymax=2000,
    legend pos=south east, legend style={font=\footnotesize, fill=white},
    legend cell align=left,
]
% power-law fits t = a |T|^b (linear regression in log-log space)
\addplot[colOX, thick, domain=15:520, samples=40, forget plot] {0.00128872*x^1.5266};
\addplot[colPMX, thick, domain=15:520, samples=40, forget plot] {0.00178166*x^1.4213};
\addplot[colOXC, thick, domain=15:520, samples=40, forget plot] {0.0962152*x^1.3598};
\addplot[colPMXC, thick, domain=15:520, samples=40, forget plot] {0.117832*x^1.3189};
% measured points
\addplot[only marks, mark=*, colOX, mark size=1.7pt] table[x=tasks,y=ox] {figures/data/cpu_time.dat};
\addplot[only marks, mark=square*, colPMX, mark size=1.6pt] table[x=tasks,y=pmx] {figures/data/cpu_time.dat};
\addplot[only marks, mark=triangle*, colOXC, mark size=2pt] table[x=tasks,y=ox_cpm] {figures/data/cpu_time.dat};
\addplot[only marks, mark=diamond*, colPMXC, mark size=2pt] table[x=tasks,y=pmx_cpm] {figures/data/cpu_time.dat};
\legend{ox,pmx,ox\_cpm,pmx\_cpm}
\end{axis}
\end{tikzpicture}

\caption{Czas CPU vs liczba zadań (skala log--log, równy budżet). Cztery
konfiguracje tworzą jedną chmurę o~zbliżonym koszcie. Markery -- pomiary, linie --
dopasowanie potęgowe $t=a\,|T|^{b}$.}
\label{fig:cpuscatter}
\end{figure}

\paragraph{Strojenie hiperparametrów.}
Bazowe parametry z~tabeli~\ref{tab:hparams} poddaliśmy strojeniu narzędziem
Optuna~\cite{akiba2019} (sampler TPE, \textit{ang. Tree-structured Parzen
Estimator}), osobno dla każdej z~czterech konfiguracji. Strojono $p_c$, $p_m$,
populację, rozmiar turnieju oraz (dla wariantów CPM) liczbę prób
\textit{pull-left}, minimalizując średnią lukę do~LB \emph{przy tym samym
budżecie $B$}. Aby uniknąć przeuczenia, strojenie prowadzono na~6~instancjach
(zbiór strojeniowy), a~porównanie strojonych parametrów z~bazowymi wykonano na~14
\emph{niewidzianych} instancjach walidacyjnych. Test H3 (sparowany Wilcoxon)
wykazał, że strojenie pomaga: w~31 z~56 testów (4 konfiguracje $\times$ 14
instancji) różnica jest istotnie na~korzyść parametrów strojonych, a~w~żadnym
na~korzyść bazowych (pozostałe istotne wyniki to remisy median). Największy zysk
odnotował najsłabszy operator OX (12~istotnych poprawek), podczas gdy czysty PMX
prawie się nie zmienił (strojenie potwierdziło, że literaturowe wartości były dla
niego niemal optymalne). Co istotne, \emph{ranking konfiguracji się nie zmienia}:
po~strojeniu nadal najlepszy jest czysty PMX (średnia luka na~zbiorze strojeniowym
$1{,}6\%$ wobec $1{,}9\%$ dla \texttt{pmx\_cpm}), a~strojony \texttt{pmx\_cpm}
dobiera wręcz małą populację ($20$) i~niewiele prób CPM ($10$) -- czyli samo
strojenie ogranicza udział lokalnego przeszukiwania.

% =====================================================================
\section{Dyskusja i wnioski}
\label{sec:wnioski}
% =====================================================================
\begin{itemize}[itemsep=3pt, topsep=3pt]
    \item \textbf{Przy uczciwym budżecie najlepszy jest czysty \texttt{pmx}.}
    Konfiguracja ta jest najlepsza lub remisuje na 19 z~20 instancji i~dorównuje
    lub pokonuje HLFET/MCP na 19~z~nich. Dodanie CPM nie poprawia już wyniku
    silnego operatora PMX.
    \item \textbf{H1 odrzucona po~wyrównaniu budżetu.} Przewaga CPM widoczna przy
    równej liczbie pokoleń była artefaktem $20$--$40\times$ większego nakładu na
    pokolenie. Przy równym budżecie ewaluacji CPM istotnie \emph{szkodzi} PMX na
    trudnych grafach (wszystkie 9 istotnych pogorszeń to pary PMX), a~pomaga
    jedynie słabszemu OX. Jest to też morał metodyczny: porównania metaheurystyk
    należy prowadzić przy równym budżecie funkcji celu.
    \item \textbf{H2 odrzucona w~stronę przeciwną.} Na grafach rzadkich nie
    dominuje OX, lecz PMX (0 z~14 istotnych testów rzadkich przemawia za~OX,
    niezależnie od~budżetu). W~tej reprezentacji i~dekoderze PMX jest
    bezpieczniejszym wyborem operatora.
    \item \textbf{Strojenie pomaga, ale nie zmienia wniosków.} Optymalizacja
    hiperparametrów (Optuna, walidacja na~zbiorze niewidzianym) istotnie poprawia
    wyniki, najbardziej dla~OX, lecz ranking pozostaje ten sam -- czysty PMX jest
    najlepszy, a~strojenie samo ogranicza rolę CPM.
    \item \textbf{Koszt jest teraz porównywalny.} Przy równym budżecie cztery
    konfiguracje mają zbliżony czas CPU, więc o~wyborze decyduje wyłącznie jakość.
    \item \textbf{Najtrudniejsza instancja: \texttt{robot\_t88}.} Tu nawet najlepszy
    AE (\texttt{pmx}, mediana $741$) przegrywa z~heurystyką HLFET/MCP ($694$) --
    metaheurystyka nie dominuje prostych heurystyk w~każdym przypadku.
    \item \textbf{Kierunki dalszych prac.} Tańsza, budżetowana wersja lokalnego
    przeszukiwania, która nie traci na~uczciwym budżecie, oraz osobny eksperyment
    z~wyłączeniem mechanizmu naprawy w~celu weryfikacji hipotezy o~przewadze PMX.
\end{itemize}

% =====================================================================
\section{Wykorzystanie dużych modeli językowych}
% =====================================================================
Przy pisaniu tej dokumentacji korzystaliśmy z~dużych modeli językowych (LLM).
Pomogły nam one przy: przejrzeniu tekstu i~zebraniu uwag, poprawkach językowych
i~gramatycznych, formatowaniu tekstu oraz wykresów (kod TikZ/pgfplots), a~także
przy zbudowaniu i~sprawdzeniu spisu literatury. Wszystkie fragmenty wygenerowane
lub zmodyfikowane w~ten sposób zostały przez nas zweryfikowane.

\begin{thebibliography}{99}
\bibitem{moscato1989} P.~Moscato, \textit{On Evolution, Search, Optimization, Genetic Algorithms and Martial Arts: Towards Memetic Algorithms}, Caltech Concurrent Computation Program, Report C3P 826, 1989.
\bibitem{graham1969} R.~L. Graham, \textit{Bounds on multiprocessing timing anomalies}, SIAM Journal on Applied Mathematics, vol.~17, no.~2, pp.~416--429, 1969.
\bibitem{davis1985} L.~Davis, \textit{Applying adaptive algorithms to epistatic domains}, w: Proc. IJCAI, pp.~162--164, 1985.
\bibitem{goldberg1985} D.~E. Goldberg, R.~Lingle, \textit{Alleles, loci, and the traveling salesman problem}, w: Proc. 1st Int. Conf. on Genetic Algorithms, pp.~154--159, 1985.
\bibitem{kelley1959} J.~E. Kelley, M.~R. Walker, \textit{Critical-path planning and scheduling}, w: Proc. Eastern Joint Computer Conference, pp.~160--173, 1959.
\bibitem{tobita2002} T.~Tobita, H.~Kasahara, \textit{A standard task graph set for fair evaluation of multiprocessor scheduling algorithms}, Journal of Scheduling, vol.~5, no.~5, pp.~379--394, 2002.
\bibitem{adam1974} T.~L. Adam, K.~M. Chandy, J.~R. Dickson, \textit{A comparison of list schedules for parallel processing systems}, Communications of the ACM, vol.~17, no.~12, pp.~685--690, 1974.
\bibitem{wu1990} M.-Y. Wu, D.~D. Gajski, \textit{Hypertool: A programming aid for message-passing systems}, IEEE TPDS, vol.~1, no.~3, pp.~330--343, 1990.
\bibitem{wilcoxon1945} F.~Wilcoxon, \textit{Individual comparisons by ranking methods}, Biometrics Bulletin, vol.~1, no.~6, pp.~80--83, 1945.
\bibitem{hooker1995} J.~N. Hooker, \textit{Testing heuristics: We have it all wrong}, Journal of Heuristics, vol.~1, no.~1, pp.~33--42, 1995.
\bibitem{deap2012} F.-A. Fortin \textit{et al.}, \textit{DEAP: Evolutionary algorithms made easy}, JMLR, vol.~13, pp.~2171--2175, 2012.
\bibitem{akiba2019} T.~Akiba, S.~Sano, T.~Yanase, T.~Ohta, M.~Koyama, \textit{Optuna: A next-generation hyperparameter optimization framework}, w: Proc. 25th ACM SIGKDD, pp.~2623--2631, 2019.
\bibitem{ullman1975} J.~D. Ullman, \textit{NP-complete scheduling problems}, J. Computer and System Sciences, vol.~10, no.~3, pp.~384--393, 1975.
\bibitem{miller1995} B.~L. Miller, D.~E. Goldberg, \textit{Genetic algorithms, tournament selection, and the effects of noise}, Complex Systems, vol.~9, no.~3, pp.~193--212, 1995.
\bibitem{whitley1994} D.~Whitley, V.~S. Gordon, K.~Mathias, \textit{Lamarckian evolution, the Baldwin effect and function optimization}, w: PPSN III, LNCS 866, Springer, pp.~6--15, 1994.
\bibitem{virtanen2020} P.~Virtanen \textit{et al.}, \textit{SciPy 1.0: Fundamental algorithms for scientific computing in Python}, Nature Methods, vol.~17, pp.~261--272, 2020.
\bibitem{mitchell2011} S.~Mitchell, M.~O'Sullivan, I.~Dunning, \textit{PuLP: A linear programming toolkit for Python}, Univ. of Auckland, 2011.
\bibitem{dejong1990} K.~A. De Jong, W.~M. Spears, \textit{An analysis of the interacting roles of population size and crossover in genetic algorithms}, w: PPSN I, LNCS 496, Springer, pp.~38--47, 1990.
\bibitem{grefenstette1986} J.~J. Grefenstette, \textit{Optimization of control parameters for genetic algorithms}, IEEE Trans. Systems, Man, and Cybernetics, vol.~16, no.~1, pp.~122--128, 1986.
\bibitem{omara2010} F.~A. Omara, M.~M. Arafa, \textit{Genetic algorithms for task scheduling problem}, Journal of Parallel and Distributed Computing, vol.~70, no.~1, pp.~13--22, 2010.
\bibitem{hou1994} E.~S.~H. Hou, N.~Ansari, H.~Ren, \textit{A genetic algorithm for multiprocessor scheduling}, IEEE Transactions on Parallel and Distributed Systems, vol.~5, no.~2, pp.~113--120, 1994.
\end{thebibliography}

\end{document}
